\subsection{无单位元代数的情形}

设 A 是 K = R 或 C 上的完备可赋范代数，不必含单位元。用 ( \(\widetilde{A}, e\) ) 表示由 A 添加单位元而得到的幺代数。它是可赋范且完备的。

设 \(x\) 是 A 的元素。若 K = C，用 \(\mathrm{Sp}'(x) = \mathrm{Sp}_{\widetilde{\mathrm{A}}} (x)\) 表示 \(x\) 相对于 \(\widetilde{\mathrm{A}}\) 的谱，并考察芽 \(f \in \mathcal{O}(\mathrm{Sp}'(x))\)。若 K = R，用 \(\mathrm{Sp}'(x)\) 表示元素 \(x\)（属于 \(\widetilde{\mathrm{A}}\)）的复谱，并考察芽 \(f \in \mathcal{O}_{\mathbf{R}}(\mathrm{Sp}'(x))\)。在这两种情形，0 都属于 \(\mathrm{Sp}'(x)\)，并且元素 \(f(x)\)（属于 \(\widetilde{\mathrm{A}}\)）属于 A 当且仅当 \(f\) 满足 \(f(0) = 0\)。事实上，投影 \(\pi: \widetilde{\mathrm{A}} \to \mathrm{Ke}\) 是核为 A 的连续态射，并且有 \(\pi(f(x)) = f(\pi(x)) = f(0)\)。

